
%\addcontentsline{toc}{section}{Г. \,Реалізація проєкту Периметр-Фортеця }

% \addtocontents{toc}{\protect\setcounter{tocdepth}{-1}}

\section {\centering Реалізація проєкту ``Кібер-Ревізор''}

У сучасному цифровому середовищі робоча станція кінцевого користувача є однією з найбільш критичних, але водночас найуразливіших точок ІТ-інфраструктури. Зростання кількості цільових атак, таких як фішинг, програми-вимагачі (Ransomware) та експлойти нульового дня, вимагає переходу від фрагментарного захисту до цілісної та багаторівневої стратегії. Робота над проєктом ``Кібер-Ревізор'' полягає у створенні такої інтегрованої моделі ``Щит для Робочої Станції''. Але не обмежується лише технічними налаштуваннями; вона поєднує організаційні принципи (політика кібергігієни), системні механізми (керування доступом та шифрування) і мережевий захист (брандмауер). Структура проєкту побудована таким чином, щоб спочатку необхідно провести ґрунтовний аналіз ризиків, а потім, базуючись на принципі ``заборона за замовчуванням'' та захисті  тріади Конфіденційність, Цілісність, Доступність, представити практичну методологію реалізації та необхідний інструментарій.

Мета такого проєкту полягає у розробці та практичній реалізації цілісної стратегії захисту кінцевого користувача, що охоплює організаційні, системні та мережеві засоби.

До проєкту можна поставити таке проблемне питання - Як забезпечити надійний та інтегрований захист конфіденційних даних на кінцевій робочій станції, враховуючи необхідність протидії гібридним загрозам (соціальна інженерія, Ransomware, мережеві атаки) за допомогою мінімального, але ефективного набору засобів?

Що стосується ключових питань, то можна запропонувати:
\begin{itemize}
    \item Як інтегрувати керування доступом (ІАА) та локальний брандмауер для забезпечення принципу найменших привілеїв?
    \item Які існують основні правові наслідки витоку персональних даних з пристрою і як їх мінімізувати
    \item Як оптимізувати стратегію резервного копіювання, враховуючи необхідність контролю цілісності даних?
\end{itemize}

Реалізація проєкту ґрунтується на трьох ключових фундаментальних засадах та комплексній стратегії:
\begin{enumerate}
    \item Захист Тріади CIA (Confidentiality, Integrity, Availability)
    \begin{itemize}
        \item Конфіденційність (C): Реалізується через повне шифрування диска (BitLocker/VeraCrypt).
        \item Цілісність (I): Забезпечується через криптографічне хешування (SHA-256) під час резервного копіювання.
        \item Доступність (A): Підтримується впровадженням стратегії резервного копіювання 3-2-1.
    \end{itemize}
    \item Принцип ``Deny by Default '' (заборона за замовчуванням) - реалізується через налаштування хостового брандмауера, який блокує весь вхідний трафік, дозволяючи лише мінімально необхідні з'єднання (наприклад, для HTTP/HTTPS та DNS).
    \item Принцип Найменших Привілеїв (PoLP) та ІАА - реалізується через розмежування облікових записів (Admin-High для адміністрування та User-Low для щоденної роботи) для мінімізації шкоди від експлойтів та Ransomware.
\end{enumerate}

Етапи виконання проєкту:
\begin{itemize}
    \item організаційна підготовка та встановлення пріоритетів;
    \item системна безпека (cia-захист);
    \item мережевий щит (deny by default);
    \item моніторинг та підтримка.
\end{itemize}

Організаційна підготовка розпочинається з початкового аудиту та ідентифікації, під час якого проводиться інвентаризація конфіденційних даних на станції та визначається, які програми потребують мережевого доступу для подальшого налаштування брандмауера. Наступний крок — вибір засобів, що передбачає встановлення рекомендованих програмних засобів (наприклад, BitLocker, Veeam Agent, Nmap) відповідно до операційної системи, що планується захищати. Кульмінацією етапу є створення облікових записів (ІАА): необхідно створити два облікові записи — Admin-High (з правами адміністратора) та User-Low (для щоденної роботи без адміністративних прав). Отже на цьому етапі обґрунтовується  реалізація принципу найменших привілеїв.

Системний захист розпочинається з шифрування диска (конфіденційність), де активується повне шифрування системного диска (BitLocker або VeraCrypt), а ключ відновлення обов'язково зберігається в безпечному місці (наприклад, на USB-носії, відокремленому від робочої станції). Далі йде налаштування резервного копіювання (доступність та цілісність), що включає щоденне інкрементне копіювання даних облікового запису User-Low за стратегією 3-2-1. У цьому процесі інтегрується механізм хешування (SHA-256), який автоматично створює хеш-суму архіву після завершення повного копіювання для контролю цілісності. Завершує етап антивірусний захист — забезпечення роботи Microsoft Defender або ClamAV у фоновому режимі з налаштуванням автоматичних оновлень сигнатур.

Мережевий захист ґрунтується на конфігурації брандмауера. Спершу встановлюється «заборона за замовчуванням» через правило ``BLOCK '' для всього вхідного трафіку (IN), а потім додаються винятки, тобто створюються дозволяючі правила (ALLOW) для необхідних служб, таких як DNS, HTTP/HTTPS, внутрішні мережеві ресурси та спеціальний SSH-доступ для адміністратора, відповідно до конфігурації. Після налаштування проводиться тестування брандмауера за допомогою Nmap (Network Mapper) з іншого пристрою в мережі для зовнішнього сканування портів. Очікуваний результат тестування полягає в тому, що Nmap повинен показати, що всі порти, крім спеціально дозволених, є заблокованими (filtered або closed).

Фінальний етап включає впровадження кібергігієни через проведення навчання користувача щодо ідентифікації фішингу (згідно з обраною політикою). Необхідний також регулярний моніторинг, щоквартально перевірятиме статус шифрування диска, актуальність антивірусних баз та успішність виконання резервних копій. Нарешті, слід проводити тестування відновлення — щорічне тестове відновлення частини даних із копії для підтвердження надійності стратегії.

У випадку виконання проєкту групою, рекомендується розподіл функціональних ролей (таблиця~\ref{tab:appendixG_role}), що забезпечує глибоке опрацювання всіх розділів модуля.

\begin{longtable}{|>{\raggedright\arraybackslash}p{3.3cm}|>{\raggedright\arraybackslash}p{4.2cm}|>{\raggedright\arraybackslash}p{7.5cm}|}

    % --- ЗАГОЛОВОК ТАБЛИЦІ ---
    \caption{Розподіл ролей у групі}
    \label{tab:appendixG_role} \\

    % --- ШАПКА ПЕРШОЇ СТОРІНКИ ---
    \hline
    \textbf{Ролі} & \textbf{Основна сфера відповідальності} & \textbf{Ключові завдання у проєкті} \\
    \hline
    \endfirsthead

    % --- ШАПКА НАСТУПНИХ СТОРІНОК (якщо таблиця переноситься) ---

\multicolumn{3}{r}{{\normalsize Продовження табл. \thetable}} \\
\hline
\textbf{Ролі} & \textbf{Основна сфера відповідальності} & \textbf{Ключові завдання у проєкті} \\
    \hline
    \endhead

    % --- НИЖНЯ ЛІНІЯ (FOOTER) ---
    \hline
    \endlastfoot

    % --- ТІЛО ТАБЛИЦІ ---

    Кібер-аналітик & 
    Аналіз ризиків та організаційний рівень &
    \begin{minipage}[t]{\linewidth}
        \vspace{1mm} % Відступ зверху
        \begin{enumerate}
            \item Розробка та обґрунтування матриці ризиків.
            \item Формулювання політики кібергігієни та аналіз правових наслідків.
            \item Складання вступу та формулювання проблемного запитання.
        \end{enumerate}
        \vspace{1mm} % Відступ знизу
    \end{minipage} \\
    \hline

    Архітектор мережі & 
    Системний захист (Тріада CIA) &
    \begin{minipage}[t]{\linewidth}
        \vspace{1mm}
        \begin{enumerate}
            \item Розробка конфігурації брандмауера за принципом ``Deny by Default''.
            \item Практичне налаштування правил брандмауера.
            \item Проведення тестування мережевого захисту за допомогою Nmap. 
        \end{enumerate}
        \vspace{1mm}
    \end{minipage} \\
    \hline

    Системний інженер & 
    Системний захист (тріада CIA) &
    \begin{minipage}[t]{\linewidth}
        \vspace{1mm}
        \begin{enumerate}
            \item Реалізація ІАА (керування доступом: admin-high vs user-low).
            \item Налаштування повного шифрування диска (конфіденційність).
            \item Розробка схеми резервного копіювання 3-2-1 з інтеграцією хешування SHA-256 (цілісність/доступність). 
        \end{enumerate}
        \vspace{1mm}
    \end{minipage} \\
    \hline

    Технічний документатор та презентатор & 
    Фінальна звітність та підготовка до захисту. &
    Перевірка узгодженості розділів, підготовка презентації. \\

\end{longtable}
При розробці проєкту учні зосереджуються на різних питаннях розділів вибіркового модуля. Перший розділ вибіркового модулю слугує основою для роботи Кібер-Аналітика. Основна увага приділяється ідентифікації загроз та аналізу їхніх джерел. Це включає розробку та обґрунтування матриці ризиків, де визначаються пріоритетні загрози, такі як комп'ютерні віруси та шкідливе програмне забезпечення (зокрема Ransomware), інтернет-шахрайство та ``крадіжка особистості''. Це також стосується необхідності формулювання політики кібергігієни як організаційного заходу протидії цим загрозам. Другий розділ вибіркового модулю стане у нагоді Системному Інженеру, який працює з найбільш технічними механізмами. Центральними питаннями є керування доступом, яке реалізовано через розмежування прав користувачів (принцип найменших привілеїв, облікові записи Admin-High та User-Low). Також ключовим є застосування криптографічних методів захисту: повне шифрування диска (BitLocker/VeraCrypt) для забезпечення конфіденційності та інтеграція хешування (SHA-256) в схему резервного копіювання для контролю цілісності даних. Учні також розглядають правові основи та відповідальність за порушення у сфері захисту інформації. Третій розділ стане у нагоді Мережевому Архітектору, який фокусується на захисті комунікацій. Ключові питання тут — це призначення та основні захисні механізми міжмережевих екранів (брандмауерів), зокрема їхня конфігурація за принципом ``Deny by Default'' для блокування зовнішніх атак, таких як сканування портів. Також важливим є використання засобів моніторингу мережевого трафіку (наприклад, Nmap або Wireshark) для тестування та підтвердження ефективності налаштувань брандмауера. До цього розділу належить і стратегія резервного копіювання (3-2-1), яка забезпечує доступність даних.

Результатом проєкту є матриця ризиків, схема політики резервного копіювання, розроблена багаторівнева модель захисту для кінцевої робочої станції, яка мінімізує найкритичніші ризики та забезпечує високий рівень захисту тріади CIA,   перелік програмних засобів, які розділені за їхнім функціональним призначенням.

Матриця ризиків є критичним аналітичним інструментом, що дозволяє визначити та кількісно оцінити потенційні загрози для робочої станції. На основі цієї оцінки (шляхом множення Ймовірності на Вплив) було встановлено пріоритети для технічних та організаційних заходів захисту, спрямовуючи ресурси на мінімізацію ризиків з найвищим рівнем.

\begin{longtable}{|
    >{\raggedright\arraybackslash}p{0.15\linewidth}| % Загроза
    >{\raggedright\arraybackslash}p{0.15\linewidth}| % Вектор
    >{\raggedright\arraybackslash}p{0.13\linewidth}| % Ймовірність
    >{\raggedright\arraybackslash}p{0.15\linewidth}| % Вплив
    >{\raggedright\arraybackslash}p{0.1\linewidth}| % Ризик
    >{\raggedright\arraybackslash}p{0.2\linewidth}| % Рекомендації
    }

    % --- ЗАГОЛОВОК ТАБЛИЦІ ---
    \caption{Матриця оцінки ризиків та загроз}
    \label{tab:risk_matrix} \\

    % --- ШАПКА ПЕРШОЇ СТОРІНКИ ---
    \hline
    \textbf{Загроза} & 
    \textbf{Вектор Атаки} & 
    \textbf{Ймовір-ність (1-5)} & 
    \textbf{Вплив (1-5)} & 
    \textbf{Рівень ризику} & 
    \textbf{Рекомендації щодо зменшення} \\
    \hline
    \endfirsthead

    % --- ШАПКА НАСТУПНИХ СТОРІНОК ---

\multicolumn{6}{r}{{\normalsize Продовження табл. \thetable}} \\
\hline
\textbf{Загроза} & 
    \textbf{Вектор Атаки} & 
    \textbf{Ймовір-ність (1-5)} & 
    \textbf{Вплив (1-5)} & 
    \textbf{Рівень ризику} & 
    \textbf{Рекомендації щодо зменшення} \\
    \hline
    \endhead

    % --- НИЖНЯ ЛІНІЯ ---
    \hline
    \endlastfoot

    % --- РЯДОК 1 ---
    Фішинг / Соціальна інженерія & 
    Електронна пошта, месенджери & 
    4 (Висока) & 
    5 (Критичний – втрата облікових даних) & 
    \textbf{20} (Високий) & 
    Регулярне навчання, двофакторна автентифікація (ІАА). \\
    \hline

    % --- РЯДОК 2 ---
    Шкідливе ПЗ & 
    Заражений файл, експлойт & 
    3 (Середня) & 
    5 (Критичний – блокування даних) & 
    \textbf{15} (Середній) & 
    Антивірусне ПЗ, ізоляція програм, регулярне резервне копіювання. \\
    \hline

    % --- РЯДОК 3 ---
    Сканування Портів & 
    Зовнішнє мережеве з'єднання & 
    2 (Низька) & 
    1 (Мінімаль-ний – інформація про порти) & 
    \textbf{2} (Низький) & 
    Хостовий брандмауер (Deny All Inbound). \\
    \hline

    % --- РЯДОК 4 ---
    Витік Даних (фізичний) & 
    Крадіжка, втрата пристрою & 
    1 (Низька) & 
    4 (Високий – втрата конфіден-ційності) & 
    \textbf{4} (Низький) & 
    Повне шифрування диска. \\

\end{longtable}


Метою схеми політики резервного копіювання є забезпечення надійної доступності даних (Availability) та їхньої цілісності (Integrity) у критичних ситуаціях, таких як вихід з ладу обладнання або атака програм-вимагачів (Ransomware), що є одним із найвищих ризиків, ідентифікованих у матриці ризиків. Для досягнення цієї мети застосовується міжнародно визнана стратегія резервного копіювання 3-2-1, яка додатково посилюється криптографічними механізмами для контролю цілісності копій.

\begin{longtable}{|
    >{\raggedright\arraybackslash}p{0.2\linewidth}| % Етап
    >{\raggedright\arraybackslash}p{0.35\linewidth}| % Дія
    >{\raggedright\arraybackslash}p{0.35\linewidth}| % Механізм
    }

    % --- ЗАГОЛОВОК ---
    \caption{Стратегія резервного копіювання та відновлення}
    \label{tab:backup_strategy} \\

    % --- ШАПКА ---
    \hline
    \textbf{Етап} & \textbf{Дія} & \textbf{Механізм Забезпечення} \\
    \hline
    \endfirsthead

    % --- ШАПКА (продовження) ---

\multicolumn{3}{r}{{\normalsize Продовження табл. \thetable}} \\
\hline
\textbf{Етап} & \textbf{Дія} & \textbf{Механізм Забезпечення} \\
    \hline
    \endhead

    % --- НИЖНЯ ЛІНІЯ ---
    \hline
    \endlastfoot

    % --- РЯДОК 1 ---
    \textbf{1. Збір} & 
    Щоденне інкрементне копіювання робочих тек і документів. & 
    \textbf{Контроль доступу (DAC):} лише User-Low має право на запис у робочу папку. \\
    \hline

    % --- РЯДОК 2 ---
    \textbf{2. Зберігання} & 
    \begin{minipage}[t]{\linewidth}
        \vspace{1mm}
        Копіювання на 3 носії:
        \begin{enumerate}
            \item Локальний зовнішній диск.
            \item Хмарний сервіс.
            \item NAS-сервер у мережі.
        \end{enumerate}
        \vspace{1mm}
    \end{minipage} & 
    \begin{minipage}[t]{\linewidth}
        \vspace{1mm}
        \textbf{Конфіденційність:} Хмарні та віддалені копії повинні бути зашифровані (AES-256).
        \vspace{1mm}
    \end{minipage} \\
    \hline

    % --- РЯДОК 3 ---
    \textbf{3. Перевірка Цілісності} & 
    Після кожного повного копіювання створюється хеш-сума (SHA-256) архіву. & 
    \begin{minipage}[t]{\linewidth}
        \vspace{1mm}
        \textbf{Цілісність:} Хеш-сума порівнюється з оригінальною для гарантії, що копія не була пошкоджена чи змінена зловмисником.
        \vspace{1mm}
    \end{minipage} \\
    \hline

    % --- РЯДОК 4 ---
    \textbf{4. Відновлення} & 
    Щорічне тестове відновлення частини даних. & 
    \textbf{Доступність:} Перевіряється, чи може система відновити дані швидко та без помилок. \\

\end{longtable}

Файл конфігурація хостового брандмауера має забезпечити принцип «заборона за замовчуванням» (Deny by Default), що критично важливо для протидії зовнішньому скануванню портів і несанкціонованому віддаленому доступу, мінімізуючи ризик, ідентифікований у матриці ризиків. Правильно налаштований брандмауер гарантує, що пристрій відкритий лише для мінімально необхідних служб, тим самим посилюючи конфіденційність і цілісність даних.


\begin{longtable}{|
    >{\raggedright\arraybackslash}p{1.5cm}| % №
    >{\raggedright\arraybackslash}p{2.5cm}| % Напрямок
    >{\raggedright\arraybackslash}p{2cm}| % Дія
    >{\raggedright\arraybackslash}p{2.5cm}| % Протокол
    >{\raggedright\arraybackslash}p{3cm}| % Локальний Порт
    >{\raggedright\arraybackslash}p{3cm}| % Віддалений Порт
    }

    % --- ЗАГОЛОВОК ТАБЛИЦІ ---
    \caption{Правила конфігурації брандмауера}
    \label{tab:firewall_rules_compact} \\

    % --- ШАПКА ПЕРШОЇ СТОРІНКИ ---
    \hline
    \textbf{№} & \textbf{Напрямок} & \textbf{Дія} & \textbf{Протокол} & \textbf{Локальний порт} & \textbf{Віддалений порт} \\
    \hline
    \endfirsthead

    % --- ШАПКА НАСТУПНИХ СТОРІНОК ---

\multicolumn{6}{r}{{\normalsize Продовження табл. \thetable}} \\
\hline
\textbf{№} & \textbf{Напрямок} & \textbf{Дія} & \textbf{Протокол} & \textbf{Локальний порт} & \textbf{Віддалений порт} \\
    \hline
    \endhead

    % --- НИЖНЯ ЛІНІЯ ---
    \hline
    \endlastfoot

    % --- РЯДКИ ПРАВИЛ ---

    % Правило 1 (Дефолтна заборона)
    \multicolumn{6}{|>{\raggedright\arraybackslash}p{\linewidth}|}{ Дефолтна заборона. Заблокувати всі спроби сканування портів та несанкціонованого доступу з мережі (принцип \textbf{Deny by Default}).} \\
    \hline
    1 &
    ВХІД (IN) &
    BLOCK &
    \texttt{TCP/UDP} &
    \texttt{ANY} &
    \texttt{ANY} \\
    \hline
    

    % Правило 2 (LAN доступ)
    \multicolumn{6}{|>{\raggedright\arraybackslash}p{\linewidth}|}{ Дозвіл для внутрішньої мережі (\texttt{LAN}) доступу до мережевих ресурсів (якщо потрібно для внутрішніх служб).} \\
    \hline2 &
    ВХІД (IN) &
    ALLOW &
    \texttt{UDP} &
    \texttt{137, 138, 445} &
    \texttt{ANY} \\
    \hline
    

    % Правило 3 (HTTP/HTTPS)
    \multicolumn{6}{|>{\raggedright\arraybackslash}p{\linewidth}|}{ Дозвіл для браузерів та програмного забезпечення доступу до Інтернету (\texttt{HTTP/HTTPS}).} \\
    \hline
    3 &
    ВИХІД (OUT) &
    ALLOW &
    \texttt{TCP} &
    \texttt{ANY} &
    \texttt{80, 443} \\
    \hline
    

    % Правило 4 (DNS)
    \multicolumn{6}{|>{\raggedright\arraybackslash}p{\linewidth}|}{Дозвіл \texttt{DNS}-запитів для вирішення доменних імен.} \\
    \hline
    4 &
    ВИХІД (OUT) &
    ALLOW &
    \texttt{UDP} &
    \texttt{ANY} &
    \texttt{53} \\
    \hline
    

    % Правило 5 (Спеціальний SSH)
    \multicolumn{6}{|>{\raggedright\arraybackslash}p{\linewidth}|}{  Дозвіл для адміністратора (з його фіксованої \texttt{IP}-адреси) для безпечного керування пристроєм.} \\
    \hline
    5 (Спец) &
    ВХІД (IN) &
    ALLOW &
    \texttt{TCP} &
    \texttt{22 (SSH)} &
    \texttt{Specific-IP} \\
    \hline
    

\end{longtable}


Для практичної реалізації проєкту ``Кібер-Ревізор'' учні можуть використовувати наступні програми та інструменти, розділені за їхнім функціональним призначенням.

Для забезпечення конфіденційності, цілісності та доступності даних на системному рівні рекомендується використовувати:
\begin{itemize}
    \item BitLocker (для Windows Pro/Enterprise) або VeraCrypt - для повного шифрування системного диска. Це критично важливо для забезпечення конфіденційності даних на пристрої у разі його фізичної втрати чи крадіжки.
    \item вбудовані засоби операційної системи, такі як User Account Control (UAC) - для  реалізації принципу найменших привілеїв шляхом розмежування прав між обліковими записами Admin-High та User-Low.
    \item Microsoft Defender (для Windows) або ClamAV (для Linux) - для захисту від шкідливого програмного забезпечення
    \item Veeam Agent (безкоштовна версія) або вбудована історії файлів Windows - для створення копій даних,
    \item командний рядок ОС (certutil у Windows або sha256sum у Linux) - для генерації та перевірки хеш-сум (SHA-256)
\end{itemize}

Для посилення мережевої безпеки та аудиту трафіку застосовуються такі інструменти:

\begin{itemize}
    \item Windows Defender Firewall або UFW (Uncomplicated Firewall) у Linux – для практичної реалізації політики «заборона за замовчуванням» (Deny by Default) та налаштування правил фільтрації (для хостового брандмауера)
    \item Nmap (Network Mapper) - дозволяє перевірити, чи дійсно налаштування блокують зовнішнє сканування портів (тестування хостового брандмауера)
    \item Wireshark – для візуалізації мережевого трафіку, що допомагає у розумінні принципів мережевих загроз (сканування портів) та роботи брандмауера
    \item OpenVPN Client або WireGuard - для налаштування захищеного каналу та демонстрації безпечного віддаленого з'єднання
\end{itemize}

Отже, проєкт ``Кібер-Ревізор'' успішно реалізує цілісну модель ``Щит для Робочої Станції'', що інтегрує керування доступом (ІАА), повне шифрування диска, контрольоване хешування резервних копій та хостовий брандмауер із принципом ``Deny by Default'', забезпечуючи високий рівень захисту тріади CIA від критичних загроз, таких як Ransomware та фішинг.

\label{last_page}